\section{Introduction}
\label{sec:introduction}

Transportation systems are increasingly required to perform effectively across several dimensions simultaneously. Decision makers must maintain economic efficiency, preserve service continuity under uncertain demand, and comply with increasingly stringent environmental regulations. Nevertheless, in many logistics and supply-chain applications, the mathematical representation of system dynamics remains essentially memoryless. Inventory, backlog, and service states are commonly updated through first-order recursions that depend only on the immediately preceding period. Although this assumption is analytically convenient, it may be too restrictive when current transportation performance is influenced by the cumulative history of past operations \cite{birge2011,shapiro2009}.

In practice, transportation and distribution systems frequently exhibit persistent memory effects. Congestion generated in previous periods may reduce current effective throughput, recovery from disruptions may extend over several periods, and product deterioration, handling fatigue, or service degradation may depend on the accumulated sequence of past operations rather than solely on the most recent state. Such phenomena cannot be represented adequately through a purely Markovian balance equation. Ignoring persistence may lead to an underestimation of the long-term cost of operational imbalances, an overestimation of the system's ability to recover instantaneously, and a distorted assessment of the optimal timing of transportation decisions. Recent studies in logistics and supply-chain control indicate that these memory effects may be operationally significant, particularly in systems affected by deterioration, prolonged disruption recovery, and adaptive control mechanisms \cite{cuong2023,haque2026}.

Fractional calculus provides a mathematically consistent framework for representing this type of long-term dependence. In particular, the Caputo fractional derivative allows the evolution of a state variable to depend on a weighted history of its previous states, with weights governed by a power-law decay kernel. This property is especially relevant in logistics systems, where the consequences of past disruptions or operational imbalances often diminish gradually rather than disappearing after a single period. Compared with conventional integer-order dynamics, fractional models can preserve richer temporal dependence while retaining a clear physical and operational interpretation. Moreover, when the fractional order approaches one, the formulation converges to the classical first-order model, thereby ensuring compatibility with standard transportation models \cite{caputo1967,podlubny1999,diethelm2010}. For finite-horizon optimization, the fractional operator can be discretized efficiently through the classical L1 approximation \cite{lin2007}.

At the same time, uncertainty and sustainability have become central concerns in transportation planning. Demand uncertainty affects shipment timing, transportation-mode selection, inventory protection, and shortage exposure. Environmental regulations introduce an additional interaction between operational and ecological decisions. Under a cap-and-trade or allowance-trading mechanism, the planner does not merely incur a fixed emissions tax. Instead, the planner must determine whether to purchase additional allowances, sell unused allowances, or modify transportation decisions to comply with the available carbon budget. Consequently, transportation planning becomes a joint economic, service, and environmental decision problem involving network utilization, inventory protection, service reliability, and emissions management \cite{bektas2011,asghari2021,homayouni2023carbon,qu2021,mirzaee2024green}.

This paper develops an integrated framework that combines three major dimensions: long-memory dynamics, stochastic demand, and carbon trading. We consider a multi-period transportation network composed of suppliers, transportation modes, and customers over a finite planning horizon. First-stage decisions determine strategic or baseline commitments, such as transportation-mode activation. Second-stage decisions adapt shipment quantities, inventory levels, shortages, and carbon-trading quantities after demand scenarios are realized. The effective inventory-service state evolves according to a Caputo fractional balance equation, whereas uncertain demand is represented through a stochastic drift-diffusion process discretized into a finite scenario set. Risk aversion is incorporated through conditional value-at-risk (CVaR), enabling the decision maker to explicitly control exposure to adverse scenarios instead of optimizing expected cost alone \cite{rockafellar2000}.

The proposed framework is motivated by an important gap in the existing literature. Sustainable transportation and carbon-aware supply-chain models have advanced considerably, particularly in routing, location, intermodal planning, and network design \cite{bektas2011,asghari2021}. Similarly, stochastic transportation models have become increasingly sophisticated in their treatment of uncertain demand and recourse decisions \cite{birge2011,shapiro2009,gbadegoye2025intermodal}. However, persistence is generally represented only indirectly, and relatively few studies incorporate fractional memory directly into the operational state equation of a transportation optimization model. Consequently, limited guidance is available regarding the joint effects of long-memory dynamics, demand uncertainty, risk aversion, and carbon regulation within a unified decision framework \cite{cuong2023,haque2026}.

The objective of this manuscript is to address this gap. More specifically, the study investigates the following research questions:

\begin{enumerate}[label=\textbf{RQ\arabic*.},leftmargin=2.3em]
    \item How can Caputo-type fractional memory be embedded into a multi-period transportation model under stochastic demand while preserving computational tractability?
    
    \item How does the fractional order influence shipment timing, shortage exposure, inventory protection, and carbon-trading decisions?
    
    \item How should risk aversion and carbon allowances be incorporated when the transportation state is governed by long-memory dynamics?
    
    \item Which computational framework is appropriate for solving the resulting stochastic mixed-integer optimization model at a moderate scenario scale?
\end{enumerate}

To answer these questions, this paper makes five main contributions. First, it introduces a fractional stochastic transportation model in which the Caputo derivative governs the evolution of an effective inventory-service state. This formulation captures the persistent influence of previous transportation and service imbalances. Second, it integrates carbon trading and CVaR-based risk aversion within the same optimization framework, thereby linking economic performance, environmental compliance, and protection against adverse demand scenarios. Third, it derives a deterministic equivalent of the fractional state dynamics using the classical L1 discretization scheme \cite{lin2007}, making the formulation suitable for finite-horizon mixed-integer optimization. Fourth, it establishes several theoretical properties of the model, including the existence of an optimal solution, the absence of simultaneous allowance buying and selling under a positive bid-ask spread, and a Lipschitz-type scenario-stability result. Fifth, it develops a solution methodology combining sample-average approximation (SAA) with progressive hedging (PH) to exploit the scenario structure of the stochastic problem \cite{rockafellar1991,shapiro2009}.

The remainder of the paper is organized as follows. Section~2 reviews the relevant literature on stochastic transportation, sustainable logistics, carbon regulation, and fractional-order modeling. Section~3 introduces the fractional-calculus and stochastic-programming concepts required for the proposed formulation. Section~4 presents the fractional stochastic transportation model. Section~5 establishes its principal theoretical properties. Section~6 describes the proposed solution methodology and parameter-calibration procedure. Section~7 reports and analyzes the computational experiments. Section~8 discusses the managerial and modeling implications of the results. Section~9 discusses the limitations of the current study, identifies future research directions and concludes the paper. All notation, scenario-generation procedures, proofs, and calibration details are self-contained within the main text.
